\begin{tikzpicture}[scale=1.15]
% 上半平面回路
\fill[gray!10] (2.2,0) arc (0:180:2.2) -- cycle;
\draw[gray!55,dashed] (0,0) circle (2.35);
\draw[axisline] (-2.7,0) -- (2.95,0) node[right,font=\small] {$\mathrm{Re}\,\omega$};
\draw[axisline] (0,-1.55) -- (0,2.35) node[above,font=\small] {$\mathrm{Im}\,\omega$};
% 路径
\draw[thick] (2.2,0) arc (0:118:2.2);
\draw[thick,->] (118:2.2) arc (118:131:2.2);
\draw[thick] (131:2.2) arc (131:180:2.2);
\draw[thick,->] (-2.2,0) -- (-0.05,0);
\draw[thick,->] (0.05,0) -- (2.2,0);
\node[font=\small] at (1.62,1.30) {$C_R$};
% 四个极点
\foreach \x/\y in {1.05/0.45, -1.05/0.45, 1.05/-0.45, -1.05/-0.45} {
  \draw[very thick] (\x-0.07,\y-0.07) -- (\x+0.07,\y+0.07);
  \draw[very thick] (\x-0.07,\y+0.07) -- (\x+0.07,\y-0.07);
}
\node[right,font=\footnotesize] at (1.13,0.48) {$-\omega_2$};
\node[left,font=\footnotesize] at (-1.13,0.48) {$-\omega_1$};
\node[right,font=\footnotesize] at (1.13,-0.48) {$\omega_1$};
\node[left,font=\footnotesize] at (-1.13,-0.48) {$\omega_2$};
% 标注
\node[font=\footnotesize,anchor=west,align=left] at (-2.55,-1.15)
  {$\omega_{1,2}=\mp\Omega-i\dfrac{\gamma}{2}$：\\
   两个在下半平面、两个在上半平面；\\
   上半圆闭合时只捕获上面两个};
\draw[aux] (-0.95,0.35) -- (-0.66,0.40);
\end{tikzpicture}
